% =========================================================================
\chapter{GraphHarnessX: A Graph-Native Harness Runtime}
\label{ch:ghx}
% =========================================================================

\claim{Execution becomes a first-class, machine-queryable object, attached
to the unmodified official loop from outside --- with the readout ladder as
the system's organizing frame and every mechanism claim carried by a live
record.}

Chapter~\ref{ch:baseline} established that the loop's evidence about its
own execution is wrong in a way that follows from how it is computed. This
chapter describes the system built in response. GHX makes the execution
itself a recorded, typed, queryable object, so that a question about what
happened during a run can be answered by reading a graph rather than by
pattern-matching over text.

The chapter is organized around a discipline rather than around a feature
list. Each section states what was built, what it is entitled to claim, and
--- where the honest answer is narrower than the obvious one --- what it is
not entitled to claim. \S\ref{sec:construction-limits} collects the
limitations that follow from the construction itself, disclosed here rather
than deferred to the threats section, because a limitation that follows
from a design decision belongs beside the design decision.

% -------------------------------------------------------------------------
\section{Design constraints}
\label{sec:constraints}
% -------------------------------------------------------------------------

Figure~\ref{fig:loop} places every mechanism of this chapter against the
stage of the official loop it attaches to. Five constraints were fixed
before any code was written. Each exists to close a specific way the
experiment could have been made meaningless.

\begin{figure}[t]
\centering
\resizebox{\textwidth}{!}{% The official six-stage loop (top row) and the seams at which GHX attaches
% (bottom row, dashed). Input from a figure environment; needs
% \usetikzlibrary{positioning,arrows.meta,fit,calc}.
\begin{tikzpicture}[
  font=\footnotesize,
  stage/.style={draw, rounded corners=2pt, align=center, minimum height=1.05cm,
                text width=2.25cm, inner sep=3pt, fill=white},
  ghx/.style={draw, dashed, rounded corners=2pt, align=center, minimum height=1.0cm,
              text width=2.25cm, inner sep=3pt, fill=gray!8},
  flow/.style={-{Latex[length=2mm]}, thick},
  seam/.style={-{Latex[length=1.6mm]}, densely dotted, thick, gray!70!black},
  node distance=0.34cm and 0.34cm]

  % ---- official loop -----------------------------------------------------
  \node[stage] (solver) {\textbf{Solver batch}\\100 tasks under\\config $C_k$};
  \node[stage, right=of solver] (dig) {\textbf{Digester}\\one dossier\\per task};
  \node[stage, right=of dig] (plan) {\textbf{Planner}\\one landscape};
  \node[stage, right=of plan] (evo) {\textbf{Evolver}\\candidate edits\\(manifest + config)};
  \node[stage, right=of evo] (crit) {\textbf{Critic}\\verdict; may\\ask the Evolver};
  \node[stage, text width=2.55cm, right=of crit] (gates) {\textbf{Five gates}\\structure, novelty,\\canonicalize,\\counterfactual, replay};
  \node[stage, right=of gates] (commit) {\textbf{Commit}\\ship $\mid$ no-op;\\rollback if\\hit rate $<0.5$};

  \draw[flow] (solver) -- (dig);
  \draw[flow] (dig) -- (plan);
  \draw[flow] (plan) -- (evo);
  \draw[flow] (evo) -- (crit);
  \draw[flow] (crit) -- (gates);
  \draw[flow] (gates) -- (commit);
  \draw[flow] (commit.north) -- ++(0,0.45) -| node[pos=0.25, above, font=\scriptsize] {$C_{k+1}$: the next round's configuration} (solver.north);

  % ---- GHX seams ---------------------------------------------------------
  \node[ghx, below=0.8cm of solver] (rec) {\textbf{Recorder}\\composition graph $G$,\\execution DAG $U$\\(observed edges only)};
  \node[ghx, below=0.8cm of $(dig)!0.5!(plan)$] (cone) {\textbf{Graph evidence}\\ancestor cones,\\pass/fail contrast,\\injection manifest};
  \node[ghx, below=0.8cm of evo] (typed) {\textbf{Typed candidates}\\seven graph ops,\\atomic commit,\\nested hashes};
  \node[ghx, below=0.8cm of gates] (six) {\textbf{Sixth gate}\\claimed node must be\\in the failure's graph;\\replay-gate record};
  \node[ghx, below=0.8cm of commit] (lad) {\textbf{Readout ladder}\\L0--L2 live;\\L3 by hand};

  \draw[seam] (solver.south) -- (rec.north);
  \draw[seam] (rec.east) -- (cone.west);
  \draw[seam] (cone.north) -- ($(dig.south east)!0.5!(plan.south west)$);
  \draw[seam] (typed.north) -- (evo.south);
  \draw[seam] (six.north) -- (gates.south);
  \draw[seam] (lad.north) -- (commit.south);

  % ---- legend, below everything ------------------------------------------
  \node[anchor=north, align=center, font=\scriptsize, text width=11cm]
    at ($(current bounding box.south)+(0,-0.15)$)
    {solid boxes: the released loop, run unmodified in the no-graph arm;\ \ dashed boxes: GHX, attached at seams,
     every flag off by default, every seam restored on exit};
\end{tikzpicture}
}
\caption[The six-stage loop and where GHX attaches]{The released
six-stage loop (solid) and the seams at which GHX attaches (dashed).
The recorder writes the composition graph and the unfolded execution
DAG while the solver batch runs; cone evidence enters the Digester's and
Planner's context through a recorded injection manifest; candidates
become typed graph edits; a sixth gate and a replay record sit beside the
five official gates; and the readout ladder names what is known about an
edit at commit time. With every flag off the dashed boxes are absent and
the loop is byte-for-byte the released one.}
\label{fig:loop}
\end{figure}

\begin{itemize}
  \item \textbf{Byte-pinned substrate.} The official system is vendored at
        a fixed revision under a sha256 manifest enforced by the test
        suite. The baseline arm therefore runs the authors' system, not our
        reading of it. Without this, every between-arm difference is
        confounded with our own maintenance.
  \item \textbf{Zero core edits.} No byte of the vendored tree is modified.
        Every GHX mechanism attaches at a seam. This is what makes the
        first constraint checkable rather than aspirational.
  \item \textbf{Seams restore in \texttt{finally}.} A crash must not leave
        a modified loop behind for the next round to inherit. Restoration
        is unconditional.
  \item \textbf{Flags default off.} Arm separation is a property of the
        configuration, not of operator discipline. A campaign flown with
        no flags set is byte-equivalent to the official system in its
        behaviour.
  \item \textbf{A gate never seen to fail is not yet a gate.} Every
        acceptance gate added here is mutation-tested: a deliberately
        malformed candidate must be refused, and the refusal recorded,
        before the gate is considered to exist.
\end{itemize}

% -------------------------------------------------------------------------
\section{The typed executable composition graph}
\label{sec:composition-graph}
% -------------------------------------------------------------------------

The harness's configuration is re-expressed as a typed graph. Six node
types cover the executable surface: the ten fixed runloop hook points,
processors, typed slots, bundles, skills, and tools. Twelve edge types
connect them --- ten declared or structural (attachment, ordering, declared
slot reads and writes, membership, mutual exclusion, and others) and two
observed at runtime (control and data).

The three strongly-typed edge families of AgentFlow's agent dependency
graph~\cite{agentflow} --- the structured-workflow system of that name, not
the unrelated RL-optimised agent sharing it --- are adopted directly; its
node vocabulary is
deliberately not, for the reason given in \S\ref{sec:structured-spaces}:
that vocabulary describes a design space, whereas this one has to describe
a specific running harness.

Two properties are enforced by test. There is a \textbf{single ordering
authority}: the graph is the only thing that decides processor order, so
the graph and the executed pipeline cannot disagree. And there is a
\textbf{parity door} against the legacy runloop: given the same inputs, the
graph-driven and the original runloop must produce byte-identical
behaviour.

% -------------------------------------------------------------------------
\section{The unfolded execution graph}
\label{sec:unfolded}
% -------------------------------------------------------------------------

The composition graph describes what the harness \emph{is}. The unfolded
execution graph $U$ records what one invocation of it \emph{did}: a
directed acyclic graph over processors, tool calls, and nested agents,
written by the runtime as execution proceeds.

Acyclicity is constructive rather than checked: every event receives a
global ordinal, and edges only ever run from lower to higher, so a cycle
cannot be represented. Data edges are computed by reaching definitions ---
the last writer of a slot to each reader of it --- and cross-checked
against the runtime's own slot provenance.

Three honesty rules govern the recorder, and they are presented as such
because they are the reason $U$ can be used as evidence at all:

\begin{enumerate}
  \item \textbf{Edges are never invented.} An edge exists in $U$ only if
        the runtime observed the dependency it represents.
  \item \textbf{An unattributable event is marked \texttt{UNGRAPHED}, not
        given a plausible identifier.} The recorder is permitted to say it
        does not know.
  \item \textbf{Cross-layer frontiers are reported, never silently
        crossed.} Where the graph's knowledge ends --- at a nested agent
        boundary, at an external call --- the boundary appears in the
        record.
\end{enumerate}

These rules cost coverage and buy admissibility. A graph that guesses is
not usable as the input to an acceptance gate, because a gate that trusts a
guessed edge is a gate that can be talked past.

% -------------------------------------------------------------------------
\section{Identity: three nested hashes}
\label{sec:identity}
% -------------------------------------------------------------------------

Three questions about a candidate get confused with one another in a loop
that edits itself through free text: is this the same proposal we saw
before, is this the same thing we actually deployed, and did it behave the
same way? GHX separates them with three nested hashes.

\begin{itemize}
  \item \textbf{Genotype} --- structural identity of the configuration,
        with runtime-only processors excluded. Two candidates with the same
        genotype are structurally the same candidate. This is the
        deduplication key.
  \item \textbf{Deployment} --- genotype plus the runtime overlay: what
        actually ran.
  \item \textbf{Phenotype} --- deployment plus the observed execution
        edges: what it actually did.
\end{itemize}

The containment $\text{genotype} \subseteq \text{deployment} \subseteq
\text{phenotype}$ is nailed by a hash-contract test. The nearest precedent
is the archive identity of the Darwin G\"odel Machine~\cite{dgm}, which keeps
lineage identity for a population of self-modifying agents; the difference
is that the layering here exists to separate \emph{declaration} from
\emph{observation}, which is the distinction a free-text edit surface
destroys.

% -------------------------------------------------------------------------
\section{Causal queries over the execution graph}
\label{sec:cones}
% -------------------------------------------------------------------------

The query the loop needs, and cannot express over flat logs, is: given this
failure, what could have produced it? In $U$ this is the \textbf{ancestor
cone} of the failing event --- everything upstream of it. The full cone
follows control and data edges; the data cone restricts to data flow alone.

The cone is small. Over the graph campaign's early failed cohort ($n{=}56$,
rounds R0--R1), the median trajectory is a 1{,}186-node, 606{,}000-character
record, and its full cone is \textbf{48 nodes}, its data cone \textbf{42}
--- about $4\%$ of the invocation \F{55}.

\emph{The modal wording is load-bearing and is kept everywhere in this
thesis}: the cone is what \emph{could have shaped} the failure, never what
caused it. A cone is a sound over-approximation of the causal ancestry
under the recorder's honesty rules. It is not a counterfactual, and
nothing in this thesis treats it as one.

% -------------------------------------------------------------------------
\section{Evidence integration}
\label{sec:evidence-integration}
% -------------------------------------------------------------------------

Failure cones and cross-task facts are injected into the Digester's and
Planner's context. Passing tasks enter as the \emph{contrast side} rather
than as dossiers of their own: passing-cone signatures feed the lift
control --- without them, lift is explicitly not computed --- along with
flip diffs recording what a task's cone gained or lost since it last
passed, and the partition line between the passing and failing
populations. The official Digester's own digestion of all outcomes is left
untouched, so nothing downstream changes shape.

The false flow column of \S\ref{sec:evidence-audit} is replaced. Across the
post-fix graph campaign, \textbf{0 of 859} dossiers carry it \F{7}, against
a reach of $86.0\%$ in the no-graph arm \F{4}. Every injection is recorded
in a manifest, so the question ``what did the role actually see'' has an
answer that does not depend on trusting the role's own account of it.

A second, independent flag, \texttt{HARNESSX\_GHX\_PLANNER\_SENSES}, extends
this seam to give the Planner fate-bucket base rates for target selection
and the Critic the same rates for pricing predicted hits, but it stayed off
in every whitelist campaign, so its effect is unmeasured.

\subsection*{A representation transform, not more information}

The obvious objection to any evidence-injection result is that the treated
arm was simply given more to read. This design forecloses it by
construction and then by measurement.

By construction, the rendered cone is a \emph{map, not a payload}: it
carries structure and identifiers, and copies no trajectory content. By
measurement, the median rendered cone is \textbf{15.5\,kB} against the
\textbf{87.4\,kB} trajectory it competes with for the same context budget
($n{=}44$, R1) --- a ratio of \textbf{0.12} \F{55}.

\subsection*{The evidence gets smaller and better-founded}

Four independent measurements in this thesis point the same way, and they
are worth stating together, because separately each looks like a detail:

\begin{itemize}
  \item the cone is $\sim$$4\%$ of the invocation's nodes
        (\S\ref{sec:cones});
  \item the rendered cone is $0.12$ of the trajectory's bytes (above);
  \item $k{=}1$ observations suffice for the loop to reach the correct drug
        family --- more trajectories buy argument quality, not new families
        \F{26}--\F{30} (\S\ref{sec:clinic});
  \item digests in the graph arms carry \textbf{3.50} verified citation
        anchors each against \textbf{1.75} in the no-graph arms, while the
        share of digests failing the production anchor validator falls
        from $11.7\%$ to $1.5\%$ \F{52} (\S\ref{sec:what-graph-did}).
\end{itemize}

Together these say something narrower and more defensible than ``the graph
helped'': \emph{structuring the evidence reduces its volume by an order of
magnitude while increasing the density and validity of its citations}. That
is a claim about the evidence, and it is measured. Whether it improves
outcomes is a separate question, answered --- in the negative, at this
bed's resolution --- in Chapter~\ref{ch:results}. Which component earns the
citation improvement is a third question, and \S\ref{sec:what-graph-did}
declines to answer it, because no arm isolates the three candidate causes.

% -------------------------------------------------------------------------
\section{Verification surfaces}
\label{sec:verification}
% -------------------------------------------------------------------------

Two surfaces use $U$ to check claims that were previously taken on trust.

\paragraph{Why a sixth gate.} The official structure gate already checks
that a candidate's evidence exists --- but as \emph{text}: a citation must
resolve to a file and a step index inside it. It cannot ask whether the
component a candidate says it is repairing was ever \emph{executed} on the
failure it cites, because nothing in the official record represents
execution as anything but prose. The whitelist record shows why the question
matters: across the three graph seeds, eighteen candidates cleared all
five official gates and were then found to name a node that never
appeared in the execution graph of the failure they cited \F{53}. A
second reason is a trap in the official ladder itself: the
vendored replay gate runs a trivial synthetic task (``Reply with exactly:
OK'', two steps) in which no new tool ever fires, so a gate that checked
tool candidates against \emph{that} execution would refuse every one of
them. The sixth gate is therefore defined over a real task's replay graph
and says so when it has none.

\paragraph{An alternative reading of the eighteen kills.} An equally
consistent reading of \F{53} is that the cone-based evidence
representation itself degrades the Evolver's targeting, rather than
revealing targeting that was already bad. The free-text no-graph arm has
no equivalent gate, so its comparable mistargeting rate is unmeasured, not
zero, and no hand audit of no-graph candidates against the same criterion
--- does the node a candidate claims to repair appear in a reconstructed
execution account of the failure it cites --- was run. The eighteen kills
show the gate firing; they do not by themselves distinguish the graph
making bad targeting visible from the graph making targeting worse.

\paragraph{A sixth acceptance gate.} The official ladder has five gates,
all of which read the candidate's prose. The sixth counts the candidate's
claimed target in the recorded failing execution: if a candidate says it is
repairing a node, that node must appear in the graph of the failure it
claims to repair. Where no replay graph exists, the gate records an
\emph{honest pass} with \texttt{checked=False} --- it passes the candidate
and records that it could not check. It never reports having checked when
it did not.

\paragraph{A replay gate with a live refusal record.} Of fifteen checks,
seven produced structural refusals of the form ``edited, but never ran'',
and no refused candidate was shipped \F{11}. The refusals are the evidence
that the gate exists in the sense of \S\ref{sec:constraints}.

\emph{Refusal correctness is unverified and is stated as such.} Each
refusal rests on a single replay draw, and no counterfactual batch was run
to establish that a refused candidate would in fact have failed. The gate
is demonstrated to fire; it is not demonstrated to fire correctly.

\paragraph{What ``checkable'' names.} Checkability is not itself one of
this chapter's measured quantities; it is a label for a bundle of
independently measured facts --- the false-flow column's dossier reach
\F{7}, the citation density and validity the graph produces \F{52}, the
sixth gate's kill count \F{53}, and the replay gate's refusal record
\F{11} --- each falsifiable on its own terms rather than as one
operationalised statistic.

% -------------------------------------------------------------------------
\section{The typed candidate surface}
\label{sec:candidate-surface}
% -------------------------------------------------------------------------

Candidates become typed graph edits, validated when authored and rendered
back into the official formats so that nothing downstream changes.

\subsection*{Stated against the right baseline}

What changed is not the reach of the edit surface. Typed candidate spaces
already exist in the literature (\S\ref{sec:structured-spaces}), and more
importantly, \emph{the official loop could already do this}. In the
no-graph baseline arm it authored new processor \emph{classes} as source
files and installed them through a \texttt{file://} target in
\texttt{config.yaml} --- among them the guard processor behind this
thesis's only above-noise repair (\S\ref{sec:one-repair}), plus a dead-end
actuator, a tool-loop breaker, an evidence gate, and a numeric verifier.
\texttt{processor} was that campaign's largest candidate bucket. Whatever
GHX contributes, it is not the ability to add a node.

What changed is how an edit is \textbf{committed and identified}. Under the
official path, the Evolver composes a manifest and a \texttt{config.yaml}
as free text, and the candidate's real edit is whatever survived its own
transcription. Under GHX, every mutation is one of seven typed operations
on the processor plane --- \texttt{insert\_node}, \texttt{remove\_node},
\texttt{replace\_same\_group}, \texttt{change\_dependency},
\texttt{swap\_subgraph}, \texttt{mutate\_inactive}, \texttt{mutate\_params}
--- committed as an atomic group through a five-stage transaction or
rejected whole, with both artifacts machine-generated on every call and a
write\,$\to$\,reload\,$\to$\,re-graph genotype round-trip verifying the
write (\texttt{harnessx/ghx/graph\_proposals.py}).

The honesty rules extend into the tool text the loop reads:
\texttt{mutate\_inactive} does not verify inactivity and says so in its own
description; underscore-prefixed metadata is refused outright rather than
silently accepted.

\subsection*{Coverage, stated honestly}

The typed surface is \emph{narrower} than the official free-text one. Seven
graph operations cover the processor plane; two config-section operations
cover the tools section; prompt edits enter as \texttt{mutate\_params} on
the prompt processor's template. Editing the harness's own source code is
out of scope entirely, and is recorded as a deviation
(Appendix~\ref{app:deviations}).

GHX \emph{types a slice} of the official action space. It does not extend
it, and the slice is the claim.

\subsection*{The runtime policy engine}

The graph also supports rule slots that the loop itself fills. The specimen
is a rule the loop authored in-run, which fired 19 times across 10 tasks
with zero collateral damage and zero conversions \F{11}\F{20}.

Read honestly: plenty of fire, no harm, no cure. It is an aliveness proof
--- the channel works end to end, from the loop's own reasoning to a live
runtime rule --- and not a success story.

% -------------------------------------------------------------------------
\section{The readout ladder}
\label{sec:ladder}
% -------------------------------------------------------------------------

The system's organizing frame, and the frame this thesis proposes for the
field, is a four-rung ladder of evidence about an edit:

\begin{itemize}
  \item[\textbf{L0}] \emph{Is it in the deployed configuration?} ---
        answered by the landing gates.
  \item[\textbf{L1}] \emph{Did it ever execute?} --- answered by the replay
        gate; live, with seven refusals \F{11}.
  \item[\textbf{L2}] \emph{Did it fire, and on what?} --- answered by fire
        and false-positive counts over $U$; produced by the loop itself
        \F{38}.
  \item[\textbf{L3}] \emph{Did firing change outcomes?} --- answered by
        paired same-window trials with firing attribution; built and
        exercised \emph{by hand} in this thesis \F{34}\F{36}, and never
        wired into the loop.
\end{itemize}

\subsection*{Offered as a standard, not as internal scaffolding}

Levels-of-evidence ladders do load-bearing work in other fields --- GRADE
in clinical medicine, SAE J3016 for driving automation --- precisely
because they force a system to state which rung its claims stand on.

Every self-evolving harness surveyed in \S\ref{sec:self-evolving} admits
candidates at \textbf{L0}: it checks that an edit landed, then selects on
aggregate score. None reports L1 or above. The ladder's contribution is
that stating the rung is now possible and cheap: given a recorded execution
graph, L1 and L2 are mechanical. L3 is where this thesis's evidence says
the remaining effort belongs, and \S\ref{sec:future} says so.

% -------------------------------------------------------------------------
\section{Limitations by construction}
\label{sec:construction-limits}
% -------------------------------------------------------------------------

Two limitations follow from the design and are disclosed here rather than
in the threats section, because they are consequences of decisions made in
this chapter.

\paragraph{Removal is not measurement.} The column that replaces the false
one saturates by construction under the sliding window: it marks
$549/549$ calls as used \F{8}. The false signal is gone, and that is
verifiable; but the replacement carries no information, so its presence
cannot be read as a measurement of anything. Chapter~\ref{ch:results} states
this at the point where the removal result is reported, not only here.

\paragraph{Recorder-gated dormancy.} With unfold recording disabled,
runtime rules are mute. This is a silent-loss channel that GHX introduces
itself: an edit can be deployed, pass every gate, and never speak, because
the recorder that would let it speak is off. It is disclosed, and it is
accounted for in the efficacy trials \F{36} --- where it turned out to
matter.
